% --- Use Case U1.1 ---
\vspace{0.5cm}
\noindent
\begin{tabularx}{\textwidth}{|>{\raggedright\arraybackslash}p{4cm}|X|}
\hline
\textbf{Use-case ID} \newline \textbf{Use-case name} & U1.1 \newline \textbf{Đăng nhập vào hệ thống xác thực tập trung (SSO)} \\ \hline
\textbf{Use-case overview} & Xác thực danh tính của thành viên nhà trường (Sinh viên, Giảng viên, Cán bộ) thông qua hệ thống HCMUT\_SSO để cho phép đăng nhập vào hệ thống bãi đỗ xe. \\ \hline
\textbf{Actors} & 1. Người dùng (Sinh viên, giảng viên) \\ 
& 2. HCMUT\_SSO \\ 
& 3. HCMUT\_DATACORE. \\ \hline
\textbf{Preconditions} & 
1. Người dùng đã có tài khoản định danh chính thức của nhà trường.\newline
2. Hệ thống đang ở màn hình đăng nhập.\newline
3. Kết nối mạng tới các máy chủ xác thực khả dụng. \\ \hline
\textbf{Trigger} & Bấm nút “Đăng nhập”. \\ \hline
\textbf{Steps} & 
1. Người dùng nhập tên đăng nhập và mật khẩu.\newline
2. Người dùng nhấn nút xác nhận đăng nhập.\newline
3. Hệ thống gửi yêu cầu xác thực tới máy chủ HCMUT\_SSO.\newline
4. HCMUT\_SSO kiểm tra thông tin và trả về mã xác thực thành công.\newline
5. Hệ thống truy vấn HCMUT\_DATACORE để lấy thông tin vai trò người dùng.\newline
6. Hệ thống đồng bộ thông tin vai trò để áp dụng chính sách giá và quyền lợi.\newline
7. Hệ thống khởi tạo phiên làm việc (session) bảo mật.\newline
8. Hệ thống chuyển hướng người dùng đến màn hình điều khiển chính. \\ \hline
\textbf{Post conditions} & Người dùng đã được xác thực, phiên làm việc được thiết lập và quyền hạn đã được phân định dựa trên vai trò từ DATACORE. \\ \hline
\textbf{Exception flow} & 
- \textbf{E1 (Thông tin không chính xác):} Hệ thống hiển thị thông báo lỗi và yêu cầu nhập lại mật khẩu.\newline 
- \textbf{E2 (Lỗi kết nối):} Thông báo máy chủ không phản hồi nếu không thể kết nối SSO hoặc DATACORE.\newline
- \textbf{E3 (Tài khoản bị khóa):} Thông báo lý do và hướng dẫn liên hệ quản trị viên. \\ \hline
\end{tabularx}

\vspace{0.5cm}

% --- Use Case U1.2 ---
\noindent
\begin{tabularx}{\textwidth}{|>{\raggedright\arraybackslash}p{4cm}|X|}
\hline
\textbf{Use-case ID} \newline \textbf{Use-case name} & U1.2 \newline \textbf{Đăng ký và kích hoạt tài khoản thành viên} \\ \hline
\textbf{Use-case overview} & Cho phép thành viên nhà trường kích hoạt tài khoản trên hệ thống bãi xe bằng cách liên kết mã định danh cá nhân với dữ liệu bãi xe thông qua HCMUT\_DATACORE. \\ \hline
\textbf{Actors}& 1. Người dùng (Sinh viên, giảng viên) \\ 
& 2. HCMUT\_SSO \\ 
& 3. HCMUT\_DATACORE. \\ \hline
\textbf{Preconditions} & 
1. Người dùng là thành viên chính thức và có dữ liệu trên HCMUT\_DATACORE.\newline
2. Người dùng sở hữu thẻ định danh còn hiệu lực.\newline
3. Kết nối với HCMUT\_DATACORE ở trạng thái sẵn sàng. \\ \hline
\textbf{Trigger} & Bấm nút “Đăng ký tài khoản”. \\ \hline
\textbf{Steps} & 
1. Đăng nhập vào HCMUT\_SSO.\newline
2. Hệ thống hiển thị biểu mẫu đăng ký (MSSV/CB, SĐT, tên đăng nhập, mật khẩu).\newline
3. Hệ thống gửi yêu cầu truy vấn thông tin tới HCMUT\_DATACORE .\newline
4. HCMUT\_DATACORE xác thực và trả về họ tên, email công vụ, vai trò.\newline
5. Hệ thống kiểm tra trùng lặp mã định danh.\newline
6. Hệ thống tự động gán chính sách giá và quyền lợi tương ứng.\newline
7. Hệ thống tạo bản ghi người dùng mới và gửi email xác nhận.\newline
8. Hệ thống thông báo thành công và chuyển về màn hình đăng nhập. \\ \hline
\textbf{Post conditions} & Tài khoản mới được tạo; vai trò và chính sách giá được thiết lập; mã thẻ định danh được liên kết thành công. \\ \hline
\textbf{Exception flow} & 
- \textbf{E1 (Mã định danh không tồn tại):} Thông báo mã số không hợp lệ hoặc không phải thành viên trường.\newline 
- \textbf{E2 (Tài khoản đã tồn tại):} Thông báo tài khoản đã được kích hoạt.\newline
- \textbf{E3 (Lỗi đồng bộ dữ liệu):} Thông báo không thể xác thực với hệ thống trường. \\ \hline
\end{tabularx}

% --- Use Case U1.3 ---
\vspace{0.5cm}
\noindent
\begin{tabularx}{\textwidth}{|>{\raggedright\arraybackslash}p{4cm}|X|}
\hline
\textbf{Use-case ID} \newline \textbf{Use-case name} & U1.3 \newline \textbf{Đổi mật khẩu người dùng} \\ \hline
\textbf{Use-case overview} & Cho phép người dùng đã đăng nhập thay đổi mật khẩu hiện tại của họ để tăng tính bảo mật cho tài khoản. \\ \hline
\textbf{Actors} & Người dùng (sinh viên, giảng viên). \\ \hline
\textbf{Preconditions} & 
1. Người dùng đã đăng nhập thành công vào hệ thống bãi xe.\newline
2. Người dùng đang ở trong giao diện "Quản lý tài khoản" hoặc "Cài đặt".\newline
3. Kết nối mạng ổn định để ghi nhận thay đổi vào cơ sở dữ liệu. \\ \hline
\textbf{Trigger} & Bấm nút “Đổi mật khẩu”. \\ \hline
\textbf{Steps} & 
1. Hệ thống hiển thị biểu mẫu đổi mật khẩu gồm 3 trường: Mật khẩu hiện tại, Mật khẩu mới, và Xác nhận mật khẩu mới.\newline
2. Người dùng nhập đầy đủ thông tin vào các trường tương ứng.\newline
3. Người dùng nhấn nút "Lưu thay đổi".\newline
4. Hệ thống kiểm tra tính chính xác của mật khẩu hiện tại bằng cách đối chiếu với dữ liệu đã lưu.\newline
5. Hệ thống kiểm tra mật khẩu mới để đảm bảo tuân thủ chính sách mật khẩu mạnh.\newline
6. Hệ thống kiểm tra xem mật khẩu mới và xác nhận mật khẩu có khớp nhau hay không.\newline
7. Hệ thống cập nhật mật khẩu mới vào cơ sở dữ liệu và ghi nhật ký thay đổi (log).\newline
8. Hệ thống hiển thị thông báo "Đổi mật khẩu thành công" và yêu cầu đăng nhập lại. \\ \hline
\textbf{Post conditions} & Mật khẩu mới được lưu thành công; mật khẩu cũ không còn hiệu lực; phiên làm việc (session) cũ có thể bị hủy. \\ \hline
\textbf{Exception flow} & 
- \textbf{E1 (Mật khẩu hiện tại không đúng):} Hệ thống thông báo "Mật khẩu hiện tại không chính xác".\newline 
- \textbf{E2 (Mật khẩu mới không đủ độ mạnh):} Hệ thống yêu cầu mật khẩu bao gồm chữ hoa, chữ thường, số và ký tự đặc biệt.\newline
- \textbf{E3 (Xác nhận mật khẩu không khớp):} Hệ thống thông báo "Mật khẩu xác nhận không trùng khớp".\newline
- \textbf{E4 (Mật khẩu mới trùng mật khẩu cũ):} Hệ thống yêu cầu sử dụng mật khẩu chưa từng dùng trước đó. \\ \hline
\end{tabularx}

\vspace{0.5cm}

% --- Use Case U1.4 ---
\noindent
\begin{tabularx}{\textwidth}{|>{\raggedright\arraybackslash}p{4cm}|X|}
\hline
\textbf{Use-case ID} \newline \textbf{Use-case name} & U1.4 \newline \textbf{Quên mật khẩu} \\ \hline
\textbf{Use-case overview} & Cung cấp cơ chế khôi phục quyền truy cập cho người dùng khi họ không nhớ mật khẩu bằng cách xác thực qua thông tin định danh của trường. \\ \hline
\textbf{Actors} & Người dùng (University Member). \\ \hline
\textbf{Preconditions} & 
1. Người dùng đã có tài khoản được kích hoạt trên hệ thống.\newline
2. Người dùng đang ở màn hình Đăng nhập.\newline
3. Hệ thống có kết nối với dịch vụ gửi email. \\ \hline
\textbf{Trigger} & Bấm nút “Quên mật khẩu”. \\ \hline
\textbf{Steps} & 
1. Hệ thống hiển thị giao diện khôi phục, yêu cầu nhập Mã số sinh viên/Cán bộ hoặc Email trường.\newline
2. Người dùng cung cấp thông tin định danh và nhấn nút "Gửi mã xác thực".\newline
3. Hệ thống tạo một mã xác thực (OTP) có thời hạn và gửi tới Email của người dùng.\newline
4. Hệ thống hiển thị màn hình yêu cầu nhập mã OTP.\newline
5. Người dùng nhập mã OTP nhận được và nhấn "Xác nhận".\newline
6. Hệ thống xác thực mã và hiển thị giao diện thiết lập mật khẩu mới.\newline
7. Người dùng nhập mật khẩu mới và nhập lại để xác nhận.\newline
8. Hệ thống thông báo "Đặt lại mật khẩu thành công" và chuyển hướng về màn hình Đăng nhập. \\ \hline
\textbf{Post conditions} & Mật khẩu tài khoản được cập nhật mới trên hệ thống xác thực trung tâm; các phiên làm việc cũ bị hủy. \\ \hline
\textbf{Exception flow} & 
- \textbf{E1 (Thông tin định danh không tồn tại):} Thông báo "Tài khoản không tồn tại".\newline 
- \textbf{E2 (Mã xác thực sai hoặc hết hạn):} Hệ thống yêu cầu nhập lại hoặc nhấn gửi lại mã mới.\newline
- \textbf{E3 (Mật khẩu mới không hợp lệ):} Hệ thống yêu cầu hiệu chỉnh lại nếu không đáp ứng tiêu chuẩn hoặc không khớp. \\ \hline
\end{tabularx}